% Compile beside math-review.sty; replace all visible insertion text for publication.
% Rename and regroup functional headings around the actual mathematical thread.
% Shared files do not determine section hierarchy; no environment/count quota applies.
\documentclass[12pt,reqno]{amsart}
\usepackage{math-review}
\title[Part II: Results (standard)]{[Review topic]\\
Part II: Results and Relations (standard)}
\author{[Author name]}
\date{[Prepared date; literature cutoff date]}
\begin{document}
\begin{abstract}
\placeholder{Describe the central answers and the deeper comparison, condition analysis and proof ideas through which this article explains the theory.}
\end{abstract}
\maketitle
\tableofcontents

\section{Introduction}\label{sec:introduction}
\placeholder{Formulate the central question and the larger distinctions that organize the known results. Explain why these distinctions, rather than chronology or an exhaustive theorem list, guide the article.}
\placeholder{Explain how Section~\ref{sec:results} establishes the main answers, how Section~\ref{sec:relations} compares their force, and why the consequences in Section~\ref{sec:consequences} describe the resulting scope. Identify where fuller explanation adds understanding.}

\section{Main results in their natural setting}\label{sec:results}
\placeholder{Develop necessary definitions and source conventions, then group results by the mathematical target. Explain the role of each group before its local variations.}
\begin{theorem}[A principal result]\label{thm:principal}
\placeholder{Give the complete verified statement and exact source version, preserving every condition and qualification shared with concise. Put a shared statement body in its canonical component.}
\end{theorem}
\placeholder{Explain why the conclusion is strong, which assumptions do real work and what the decisive proof idea supplies. Deepen these explanations, not merely the number of results.}

\subsection{An important variation or alternative formulation}\label{subsec:variation}
\placeholder{Develop a useful variation with its exact range and explain its relation to the principal result. Use a separate canonical node for a genuinely different statement, not for an editorial paraphrase.}

\section{Comparisons that organize the theory}\label{sec:relations}
\placeholder{Justify the central implications, differences, equivalences or improvements. Show how convention changes are translated and where conclusions do not compare. Explain a sourced historical advance when it changes this mathematical picture.}

\section{Consequences and the reach of the answers}\label{sec:consequences}
\placeholder{Derive significant consequences and explain the remaining boundaries of the selected results. Include supporting arguments when they improve understanding. Add supplementary families only for a clear mathematical role, rather than restoring every omitted concise branch.}

\templatereferences
\end{document}
